%==============================================================================
%  CHAPTER 5 -- RESULTS AND DISCUSSION
%
%  The chapter the examiner reads most carefully, and the one that earns the mark.
%
%  RESULTS    = what the data show. Neutral, factual, with uncertainty.
%  DISCUSSION = what it means, why it happened, how it compares with the
%               literature, and where it does not hold.
%
%  You may interleave them section by section (recommended) or keep them
%  separate. Either way the reader must always know which one they are reading.
%
%  Rules that matter:
%   - Every measured quantity carries an uncertainty. A number without one is
%     an opinion.
%   - Every figure is referred to in the text (Figure~\ref{fig:x}) and is
%     discussed. A figure nobody mentions should be deleted.
%   - Axis labels carry units, and the font must be readable at print size.
%   - Report what disagreed with your expectation. That is evidence, not failure.
%==============================================================================
\chapter{Results and discussion}
\label{ch:results}

%------------------------------------------------------------------------------
\section{Preliminary characterisation}
\label{sec:results-preliminary}
%  Whatever had to work before the main result could be trusted: repeatability
%  checks, noise floor, calibration verification, sensitivity analysis.

%------------------------------------------------------------------------------
\section{Main results}
\label{sec:results-main}
%
%  \begin{figure}[H]
%    \centering
%    \includegraphics[width=0.8\textwidth]{Figures/example}
%    \caption{A caption that can be understood without reading the body text:
%             what is plotted, under what conditions, and what to notice.}
%    \label{fig:example}
%  \end{figure}

%------------------------------------------------------------------------------
\section{Discussion}
\label{sec:discussion}
%  Answer, explicitly, the question posed in Section~\ref{sec:research-question}.
%  Compare quantitatively with the works cited in Chapter 1, and explain the
%  mechanisms behind what you observed.

%------------------------------------------------------------------------------
\section{Limitations}
\label{sec:limitations}
%  The conditions under which your conclusions do NOT hold, and why.
%
%  Students are reluctant to write this section. Do not be: stating your own
%  limits shows that you understand your work, and it stops an examiner from
%  raising them as though you had not noticed.
